\chapter{Jock Itch (Tinea Cruris)}

\section*{Definition}
Jock itch (tinea cruris) is a common fungal infection of the groin and inner thighs caused by dermatophytes. It is more common in men and in athletes, especially in warm, humid conditions.

\section*{What Happens in Your Body}
Dermatophyte fungi (most commonly Trichophyton rubrum) invade the superficial keratin layer of the skin. The warm, moist environment of the groin area promotes fungal growth. The infection elicits an inflammatory response that causes the characteristic red, itchy, ring-shaped rash. The scrotum is typically spared, which helps distinguish it from other conditions.

\section*{Causes}
\begin{itemize}
\item Excessive sweating and moisture in the groin area
\item Tight, non-breathable clothing
\item Obesity (increased skin folds and moisture)
\item Immunosuppression (diabetes, HIV, corticosteroid use)
\item Sharing contaminated towels or clothing
\item Poor hygiene
\item Prior fungal infection of the feet (athlete's foot)
\item Hot, humid climate
\end{itemize}

\section*{When to See a Doctor}
\begin{itemize}
\item Intense itching in the groin and inner thigh area
\item Red, raised, scaly rash with a well-defined border
\item Rash may have central clearing (ringworm appearance)
\item Burning or stinging sensation
\item Flaking, peeling, or cracking skin
\item Often spares the scrotum
\item May extend to the buttocks and upper thighs
\end{itemize}

\section*{Related Symptoms}
\begin{itemize}
\item Athlete's foot (tinea pedis)
\item Ringworm (tinea corporis)
\item Intertrigo (moisture-associated rash)
\item Candidal intertrigo (yeast infection)
\item Contact dermatitis of the groin
\item Psoriasis
\end{itemize}

\section*{Self-Care / Home Management}
\begin{itemize}
\item Keep the groin area clean and dry
\item Apply OTC antifungal creams (clotrimazole, miconazole, terbinafine)
\item Wear loose-fitting, breathable cotton underwear
\item Change clothes after exercise
\item Use antifungal powder to reduce moisture
\item Avoid sharing towels
\item Wash underwear in hot water
\item Dry the groin area thoroughly after bathing
\end{itemize}

\section*{How Your Doctor Will Evaluate You}
\begin{itemize}
\item Clinical examination: characteristic appearance and location
\item Skin scraping with KOH preparation for microscopic fungal elements
\item Fungal culture if diagnosis is uncertain
\end{itemize}

\section*{Treatment Options}
\begin{itemize}
\item Topical antifungal creams for 2-4 weeks (terbinafine, clotrimazole, miconazole)
\item Oral antifungal medication (terbinafine, itraconazole, fluconazole) for extensive or resistant cases
\item Topical corticosteroids for inflammation (short-term)
\item Keep area dry and avoid recurrence triggers
\end{itemize}

\section*{References}
\begin{thebibliography}{99}
\bibitem{ref1} Gupta AK, Chaudhry M, Elewski B. "Tinea corporis, tinea cruris, tinea pedis, and tinea manuum." \textit{Dermatologic Clinics}, 2003;21(3):395--400.
\bibitem{ref2} Havlickova B, Czaika VA, Friedrich M. "Epidemiological trends in skin mycoses worldwide." \textit{Mycoses}, 2008;51(Suppl 4):2--15.
\bibitem{ref3} Elewski BE. "Treatment of tinea pedis and tinea cruris." \textit{Journal of the American Academy of Dermatology}, 1996;35(3):S16--S19.
\bibitem{ref4} Moriarty B, Hay R, Morris-Jones R. "The diagnosis and management of tinea." \textit{British Medical Journal}, 2012;345:e4380.
\bibitem{ref5} Blatny R, Michalek J, Huntenburg K. "The diagnosis and management of dermatophyte infections." \textit{Hautarzt}, 2014;65(11):965--974.
\bibitem{ref6} Kelly BP. "Superficial fungal infections in athletes." \textit{Sports Medicine}, 2013;43(3):199--210.

\end{thebibliography}

\clearpage
